\chapter{Method and Progress}\label{ch:method}

This chapter follows the same grouped MAPPO and centralized-training,
decentralized-execution foundation used in the previous report. The new work is
concentrated in the actor communication block, the construction of building
groups, and the selection of actor experts. This stage also investigates which
reward-function weight setting works best. The basic reward form is retained,
but several combinations of its weights are compared to study the balance
between load tracking and indoor comfort.


\section{Reward Function and Parameter Comparison}

At time step $t$, let $y_t$ be the total electricity load of all $N$
buildings, let $r_t$ be the target load, and let $e_t=y_t-r_t$ be the tracking
error. The two running load-tracking terms used by the reward are
\begin{align}
B_t &= \left|\frac{\sum_{\tau=1}^{t}e_\tau}
{\sum_{\tau=1}^{t}r_\tau}\right|,\\
E_t &= \frac{\sqrt{\frac{1}{t}\sum_{\tau=1}^{t}e_\tau^2}}
{\frac{1}{t}\sum_{\tau=1}^{t}r_\tau},
\end{align}
where $B_t$ is the running absolute NMBE term and $E_t$ is the running
CV-RMSE term. NMBE shows whether the total load is generally above or below the
target over time: a positive value means that it is generally too high, while a
negative value means that it is generally too low. The reward uses its absolute
value, so a result closer to zero means less overall bias. CV-RMSE shows
the typical size of the hour-by-hour tracking errors relative to the average
target load, and larger errors have a stronger effect on it. A lower CV-RMSE
therefore means that the load follows the target more closely. For building $i$,
the number of degrees outside the seasonal comfort range $[L_t,U_t]$ and the
corresponding violation indicator are
\begin{align}
d_{i,t} &= \max(L_t-T_{i,t},0)+\max(T_{i,t}-U_t,0),\\
I_{i,t} &= \mathbb{I}[d_{i,t}>0].
\end{align}
The complete per-building reward is therefore
\begin{equation}
\label{eq:report2-reward}
R_{i,t}=-\left[
\frac{w_{\mathrm{nmbe}}B_t+w_{\mathrm{cv\_rmse}}E_t}{N}
+w_{\mathrm{comfort}}
\left(w_{\mathrm{binary}}I_{i,t}+w_{\mathrm{degree}}d_{i,t}\right)
\right].
\end{equation}
The first two penalties are calculated from the total load and divided equally
among the $N$ buildings. The final term is calculated separately for each
building. It gives one penalty when the indoor temperature leaves the comfort
range and an additional penalty according to how many degrees it exceeds the
range. In the reported experiments, Vermont uses January for training and
February for testing, with a comfort range of $20$--$24^{\circ}$C. Texas uses
August for training and September for testing, with a comfort range of
$22$--$26^{\circ}$C.

The five commonly used parameters are the NMBE weight
$w_{\mathrm{nmbe}}$, CV-RMSE weight $w_{\mathrm{cv\_rmse}}$, overall comfort weight
$w_{\mathrm{comfort}}$, fixed comfort-violation weight
$w_{\mathrm{binary}}$, and temperature-exceedance weight
$w_{\mathrm{degree}}$. In this order, the usual baseline setting is
$(1.0,1.0,0.8,1.3,0.3)$. These objectives can conflict. For example, moving
battery or HVAC demand to improve indoor comfort can move the portfolio load
away from its target, while concentrating only on load tracking can leave less
freedom to correct indoor-temperature violations. NMBE and CV-RMSE can also
favour different behaviours because one measures overall bias and the other
measures hour-by-hour tracking error.

Four reward settings were compared: the baseline above; a stricter-comfort
setting $(1.0,1.0,1.5,3.0,1.0)$; a degree-hours follow-up
$(1.0,1.0,1.5,3.0,1.5)$; and a load-tracking-focused setting
$(5.0,5.0,1.0,1.0,1.5)$. The stricter-comfort setting keeps both load-tracking
weights at 1.0 while increasing the three comfort-related weights. The next
comparison keeps the first four weights unchanged and increases only
$w_{\mathrm{degree}}$ from 1.0 to 1.5. This controlled change tests whether a
larger penalty for the size of a temperature violation actually reduces
degree-hours. The load-tracking-focused setting then tests the opposite
direction by increasing the NMBE and CV-RMSE weights. The best setting is chosen
from the measured load-tracking and comfort results rather than from reward
magnitude alone.

\section{TarMAC Hybrid Communication}

The original implementation follows targeted multi-agent communication
(TarMAC) \cite{tarmac2019}. For hidden state $h_i\in\mathbb{R}^{256}$, it forms
three representations. The query describes what information building $i$ is
looking for. The key describes what information each building can provide and
is used for matching. The value contains the actual information to be passed.
They are calculated as
\begin{align}
q_i &= W_q h_i, & k_i &= W_k h_i, & v_i &= W_v h_i,\\
a_{ij} &= \operatorname{softmax}_j\left(\frac{q_i^\top k_j}{\sqrt{d_k}}\right),
& c_i &= \sum_j a_{ij}v_j,
\end{align}
The model compares the query of one building with the keys of all buildings.
A better match receives a larger weight, and the values are combined according
to these weights to form the received communication information $c_i$. The
division by $\sqrt{d_k}$ only prevents the matching scores from becoming too
large. The original implementation then computes an update directly from the
complete local state and $c_i$. Because the full 256-dimensional local state
enters the fusion layer without compression, repeated local information can
overwhelm the 64-dimensional communication information in this 25-building task.

TarMAC Hybrid changes the fusion path to
\begin{align}
\ell_i &= \operatorname{ReLU}(W_{\ell}h_i), \qquad
\ell_i\in\mathbb{R}^{64},\\
\Delta_i &= \tanh\!\left(W_2\tanh\left(W_1[\ell_i;c_i]
+b_1\right)+b_2\right),\\
h'_i &= h_i + \Delta_i.
\end{align}
Compared with the original TarMAC communication method, TarMAC Hybrid makes
three main changes. First, a learned $256\to64$ projection
turns the local state into a more compact representation and reduces repeated
information. Second, ReLU provides nonlinear filtering so that useful patterns
in the current input are more clearly represented. Third, a residual connection
adds the communication result to the original state instead of replacing the
original state completely. Together, these changes first simplify each
building's local information and then combine it with information from other
buildings without losing the original input. The default ReLU experiment uses
32-dimensional queries and keys and 64-dimensional values. It uses one global
communication round, meaning that all 25 buildings exchange information once
before the action is produced, and each building can consider information from
every other building rather than only buildings in its fixed group.

To determine which part of TarMAC Hybrid produces the improvement, the
experiments also compare two alternatives with the default ReLU fusion: Linear
and Gated. They separate the effect of simple dimensional reduction from the
effect of controlling how much communication information is added.

Linear fusion compresses the local state from 256 dimensions to 64 dimensions
without applying ReLU, and then combines this result directly with the received
communication information. It is used to test whether dimensional reduction
alone is sufficient.

Gated fusion learns a set of weights between zero and one that act like
adjustable switches. Based on the local state and the received information,
these switches decide how much of each part of the communication information
should be added to the local state. It is used to test whether controlling the
amount of incoming information is more effective than direct fusion.


\section{Grouping Method Settings}

As explained in Chapter~\ref{ch:background}, this task has no correct actor-group
labels, so it requires unsupervised grouping. The grouping-method ablation
compares K-means as the baseline with three alternatives: GMM, agglomerative
clustering, and balanced spectral clustering. GMM was included to allow groups
with different spreads and some overlap; agglomerative clustering was included
to remove random initialization; and balanced spectral clustering was included
to test whether directly balancing group size could improve the result.

The grouping-method ablation used one common 5F configuration to compare K-means,
GMM, agglomerative clustering, and balanced spectral clustering. All four methods
used exactly the same feature input, TarMAC Hybrid fusion, random seed, grouped
MAPPO training procedure, and evaluation procedure. This controlled design keeps
the grouping method as the main difference between the four experiments.

The grouping-method ablation uses the feature matrix
$X\in\mathbb{R}^{25\times5}$. Each row represents one building, and the five
columns are battery capacity, total HVAC rated power, mean heating demand, mean
non-shiftable load, and mean lower-comfort-bound exceedance. The same combination
is supplied to all four methods for a fair comparison. In the feature ablations, the matrix
has the more general form $X\in\mathbb{R}^{25\times d}$, where $d$ changes with
the selected feature columns. Once a feature combination is fixed for one
comparison, every run uses exactly the same data-derived matrix; the matrix
values are not randomly generated.

\noindent\textbf{K-means.}
For both four and five candidate groups, K-means assigns each building to its
nearest centre and updates the centres repeatedly. Each fit tries ten initial
centre settings to reduce the effect of an unsuitable random start. The best
balanced candidate is retained.

\noindent\textbf{GMM.}
GMM estimates the probability that each building belongs to each group and
assigns it to the group with the highest probability. It is tested with both four
and five groups. Each fit uses three starting points and allows each group to
have its own width and direction; this is the purpose of the full covariance
matrix. A small $10^{-6}$ adjustment is added to prevent unstable numerical
calculations.

\noindent\textbf{Agglomerative clustering.}
This method starts with each building in its own group and repeatedly merges the
two groups that are most similar. Ward's rule selects the merge that keeps the
buildings inside the new group as similar as possible. The process stops when
four or five groups remain. It does not use a random starting point, so the same
feature matrix produces the same result.

\noindent\textbf{Balanced spectral clustering.}
This method first gives a stronger connection to buildings with more similar
features, and then uses the connection pattern to form groups. The Hungarian
assignment used by balanced spectral reserves a fixed number of places for each
group. It is simply a matching method: it considers the possible
building-to-place assignments and finds the one with the smallest total
assignment difference. This assigns every building while keeping group sizes
equal or different by at most one. If two candidates are equally balanced, the
lower normalized-cut cost is preferred because it means that strongly connected
buildings are kept together more successfully.

Before any method is run, every selected feature column is standardized. The
standardization and the basic group-balance score are
\begin{align}
X_{ij} &= \frac{f_{ij}-\mu_j}{\sigma_j},\\
S_{\mathrm{balance}} &= \frac{\min_k n_k}{\max_k n_k},
\end{align}
where $f_{ij}$ is the original value of feature $j$ for building $i$, and
$\mu_j$ and $\sigma_j$ are the mean and standard deviation of that feature over
the 25 buildings. Standardization stops a feature with large numerical units
from dominating the grouping. In the balance score, $n_k$ is the number of
buildings in group $k$; a value closer to one means that the group sizes are more
similar. The experiments treat candidate groupings containing an empty group or
a group with only one building separately. If a candidate contains an empty
group, its selection score is set directly to $-1$; if its smallest group contains
only one building, $0.5$ is subtracted from the ratio above. The selection score
is a value used to compare whether buildings are distributed evenly across the
different candidate groupings. Normally, it is the balance score above, namely
the smallest group size divided by the largest group size. Finally, the candidate
with the highest selection score is retained.


\section{Grouping Feature Design}

Based on the results of the preceding grouping-method experiment, agglomerative
clustering performed well and was fixed as the grouping method for the subsequent
feature experiments. This keeps the clustering algorithm unchanged while the
selected feature columns are varied.

The initial experiment considered 37 possible features: seven static building
features and 30 operational features. The static features describe battery and
HVAC size, photovoltaic system size and roof area, and whether different thermal
storage systems are installed. The operational features describe non-shiftable
load, cooling, heating, domestic hot water, solar generation, occupancy, indoor
temperature, and comfort-bound exceedance using statistics such as the mean,
standard deviation, maximum, and sum.

Before the compact ablations, these 37 features were screened to nine. Three
principles were used. First, the retained features should be directly related to
the available control flexibility, electrical demand, thermal demand, or indoor
comfort. Second, highly repetitive descriptions of the same signal were reduced:
mean and maximum values were kept to represent typical and peak demand, while
several standard-deviation and sum features were removed. Third, the selection
matched the Vermont winter experiment, so heating demand and lower-comfort
exceedance were more informative than cooling demand and upper-comfort
exceedance. PV, solar, storage-presence flags, and other secondary descriptors
were not included in the compact pool because they were less direct for this
specific control question or repeated information already represented by the
retained demand and capacity variables. This screening produced the following
nine control-profile features:
\begin{itemize}
  \item \texttt{bes\_capacity\_kwh}: battery energy capacity, representing how much
  electrical energy the building can shift over time;
  \item \texttt{hvac\_total\_kw}: total rated HVAC power, representing the installed
  ability to change heating or cooling electricity use;
  \item \texttt{heating\_mean}: average heating demand, representing the building's
  typical thermal load;
  \item \texttt{heating\_max}: maximum heating demand, representing its peak thermal
  load;
  \item \texttt{nsl\_mean}: average non-shiftable load, representing typical
  background electricity demand;
  \item \texttt{nsl\_max}: maximum non-shiftable load, representing peak background
  demand;
  \item \texttt{occupants\_mean}: average occupancy, representing typical building
  use intensity;
  \item \texttt{indoor\_temp\_std}: standard deviation of indoor temperature,
  representing how much indoor temperature varies;
  \item \texttt{comfort\_lower\_excess\_mean}: average amount by which indoor
  temperature falls below the lower comfort limit, representing under-heating
  behaviour.
\end{itemize}

The full 9F control-profile description consists of all nine variables above.
The experiment first used all nine features in an agglomerative-clustering
grouped TarMAC Hybrid MAPPO experiment. The input features were then reduced
progressively:
\begin{itemize}
  \item 5F: battery capacity, HVAC rated power, mean heating demand, mean
  non-shiftable load, and mean lower-comfort-bound exceedance;
  \item 4F: the first four features above, with the comfort feature removed;
  \item 3F: battery capacity, mean heating demand, and mean non-shiftable load,
  with HVAC rated power also removed.
\end{itemize}
This reduction tests whether adding more building descriptions provides useful
grouping information or introduces redundant information.

After the feature-count ablation, the experimental results showed that the 3F
settings generally performed better. A detailed analysis is provided in
Chapter~\ref{ch:evaluation}.

Choosing any three variables from the nine-feature pool gives
$\binom{9}{3}=84$ possible combinations. Training every combination with several
random seeds would be expensive, so nine representative combinations were
selected using physical reasoning and controlled substitutions:
\begin{itemize}
  \item A: capacity + heating mean + NSL mean, covering storage flexibility,
  typical thermal demand, and typical background load;
  \item B: capacity + HVAC rating + NSL mean, testing installed HVAC capability
  instead of measured heating demand;
  \item C: capacity + heating maximum + NSL mean, testing peak rather than average
  heating demand;
  \item D: capacity + heating mean + NSL maximum, testing peak rather than average
  non-shiftable load;
  \item E: capacity + HVAC rating + NSL maximum, combining installed HVAC
  capability with peak background demand;
  \item F: capacity + heating maximum + NSL maximum, combining the two peak-demand
  descriptions;
  \item G: capacity + heating mean + lower-comfort exceedance, testing whether a
  direct comfort-related feature is more useful than non-shiftable load;
  \item H: capacity + heating mean + indoor-temperature variation, testing whether
  indoor thermal stability helps to define the groups;
  \item I: heating mean + NSL mean + lower-comfort exceedance, removing capacity to
  test whether battery size is essential to the grouping.
\end{itemize}
These combinations cover the main questions of installed capability versus
observed demand, mean versus peak demand, electrical load versus comfort and
temperature information, and grouping with or without battery capacity. Average
occupancy was not included in the nine compact combinations because its effect
is partly reflected in heating and non-shiftable load, while those two variables
are more directly related to the control task. The selected combinations are
therefore a focused screening of the most plausible choices rather than an
exhaustive search of all 84 possibilities.

For Texas, the same feature roles are retained but
\texttt{heating\_mean} is replaced with \texttt{cooling\_mean}. Thus the
cross-climate experiment tests a transferable principle---capacity, controllable
thermal burden, and background load---rather than blindly copying a Vermont
heating feature into a cooling-dominated season.


\section{Soft Router for Dynamic Actor Selection}

Fixed grouping assigns each building to one actor before training, and that actor
does not change when the building's battery level, load, or heating demand
changes. Its performance therefore depends strongly on the quality of the
initial grouping. Soft Router was introduced to address this limitation. Its
purpose is to let each building choose the most suitable actor for its current
condition, while still making use of the group actors that have already been
trained. It therefore tests whether actor selection can become more flexible and
less dependent on one fixed grouping.

The Soft Router itself is a small shared neural network: for each building, it
reads the building's current information and outputs one selection probability
for every candidate actor. Each trained group actor is one such candidate, also
called an expert. To explain how the router chooses among them, first consider
how one expert produces an action. At every time step, a building provides an observation
$o_{it}$, such as its battery SOC, current load, heating demand, indoor
temperature, and time information. The expert's \emph{encoder} turns these raw
values into a short internal summary. TarMAC adds useful information from the
other buildings, and the \emph{action head} finally converts the communicated
summary into a control action.

The encoder and action head of each pretrained expert were learned together.
Some early routers switched only the action head, so a head could receive an
internal summary produced by an encoder with which it had never been trained.
The stable version instead switches the matched encoder and action head together;
this is called \emph{full-expert routing}. It preserves the complete control
behaviour of each pretrained actor and is therefore easier to keep stable. The
following explanation focuses on this version, while the other router designs
are summarized later.

To choose a full expert, the router receives an input with two complementary parts.
First, the encoder of the actor originally assigned to each building converts
its observation into a 256-dimensional latent feature. TarMAC then communicates
among all 25 buildings. Thus, the communicated feature for building $i$ can be
written as
\begin{equation}
h_{it}=\operatorname{TarMAC}_i\!\left(
\left\{\operatorname{Enc}_{c_j}(o_{jt})\right\}_{j=1}^{N}
\right)\in\mathbb{R}^{256},
\end{equation}
where $c_j$ is the original fixed group of building $j$. Although $h_{it}$ is
associated with building $i$, it contains information received from the other
buildings. It is a learned summary rather than a directly measured physical
quantity.

The learned summary is rich, but an encoder trained mainly for action generation
may compress or obscure physical quantities that are especially important for
choosing between experts. The capacity-aware router therefore appends seven
explicit features for the same building $i$,
\begin{equation}
q_{it}=\left[
\bar C_i,\;\bar H_i,\;s_{it},\;
\bar C_i(1-s_{it}),\;\bar C_i s_{it},\;
\ell_{it},\;u_{it}
\right]\in\mathbb{R}^{7},
\end{equation}
where $\bar C_i$ and $\bar H_i$ are its normalized battery capacity and HVAC
rating, $s_{it}$ is its SOC, $\ell_{it}$ is its current non-shiftable load, and
$u_{it}$ is its current heating demand. The fourth and fifth entries explicitly
represent remaining charging space and available discharge energy. These seven
values belong only to building $i$; they are not a concatenation of the seven
features of all buildings. Cross-building information is already supplied by
$h_{it}$.

The complete router input is therefore
\begin{equation}
x^{\mathrm{router}}_{it}=[h_{it};q_{it}]\in\mathbb{R}^{263}.
\end{equation}
There is one shared router network, not one separately parameterized router per
building. The same parameters $\theta$ are applied to all 25 inputs, while the
different building states produce 25 different routing distributions. In the
reported stable full-expert run, this feed-forward network has the form
$263\rightarrow64\rightarrow5$: a linear layer, ReLU, and a second linear layer
that produces one logit for each of the five experts. For building $i$ at time
$t$, the logits are converted into probabilities,
\begin{equation}
d_{it}=\operatorname{softmax}\!\left(
\operatorname{Router}_{\theta}(x^{\mathrm{router}}_{it})/\tau
\right),
\end{equation}
where every element of $d_{it}$ is the probability of selecting one actor. The
probabilities add up to one. The temperature $\tau$ controls how decisive the
choice is: a lower value makes the router concentrate more strongly on one actor,
whereas a higher value spreads the probabilities more evenly.

At the beginning of training, immediately abandoning the original grouping can
destroy the useful behaviour already learned by the actors. The router therefore
combines its new probabilities with the original fixed assignment,
\begin{equation}
g_{it}=(1-\alpha)p_i+\alpha d_{it}.
\end{equation}
Here, $p_i$ marks the actor originally assigned to building $i$, and $\alpha$
controls how much freedom is given to the router. When $\alpha=0$, the original
fixed actor is used. As $\alpha$ increases, the router is allowed to choose other
actors more often. The experiments pass $\tau$, the starting and final values of
$\alpha$, the warm-up length, the router learning rate, and the balance strength
as training settings. This makes it possible to change how quickly and how
strongly dynamic selection is introduced without changing the building data.

During training, an expert is selected according to $g_{it}$ and generates the
building action. To protect the useful behaviour in the pretrained actors, the
stable experiment uses staged training: it first freezes all experts and updates
only the router. This raises an important question: if the experts no longer
change, how can a policy loss still train the router? The reason is that the
complete policy includes both the expert action distributions and the router's
selection probabilities. Let $\pi_{\phi_k}(a_{it}\mid s_t)$ denote the action
distribution of expert $k$, with expert parameters $\phi_k$. The complete routed
policy is the mixture
\begin{equation}
\pi_{\theta,\phi}(a_{it}\mid s_t)
=\sum_{k=1}^{K}g_{itk}(\theta)\,
\pi_{\phi_k}(a_{it}\mid s_t).
\label{eq:soft-router-mixture-policy}
\end{equation}
Freezing an expert fixes $\phi_k$ and its action distribution, but it does not
fix the complete policy. Changing the router parameters $\theta$ changes the
mixture weights $g_{itk}$ and therefore changes the probability of the executed
action. Equation~\eqref{eq:soft-router-mixture-policy} is therefore the link
between staged freezing and the PPO loss derived below.

CityLearn executes the actions and returns the same reward defined in
Equation~\eqref{eq:report2-reward}; the router has no separate reward. The shared
part measures portfolio load tracking, and the building-specific part measures
indoor comfort. After one rollout, MAPPO uses the rewards and the critic to form
an advantage estimate, which can be summarized as
\begin{equation}
\hat A_{it}=\hat R_{it}-V_{\psi}(s_t).
\end{equation}
A positive advantage means that the sampled action led to a better result than
the critic expected, while a negative advantage means that it was worse. For the
stored action, PPO recomputes the routed-policy probability and forms the ratio
\begin{equation}
r_{it}(\theta,\phi)=
\frac{\pi_{\theta,\phi}(a_{it}\mid s_t)}
{\pi_{\mathrm{old}}(a_{it}\mid s_t)}.
\end{equation}
The clipped policy loss is
\begin{equation}
\mathcal{L}_{\mathrm{PPO}}
=-\mathbb{E}_{i,t}\!\left[
\min\!\left(
r_{it}\hat A_{it},
\operatorname{clip}(r_{it},1-\epsilon,1+\epsilon)\hat A_{it}
\right)
\right].
\label{eq:soft-router-ppo-loss}
\end{equation}
Even when $\phi_k$ is frozen, Equation~\eqref{eq:soft-router-mixture-policy}
shows that $r_{it}$ and the PPO loss still depend on $\theta$. The gradient path is
\begin{equation}
\frac{\partial\mathcal{L}_{\mathrm{PPO}}}{\partial\theta}
=\frac{\partial\mathcal{L}_{\mathrm{PPO}}}{\partial\pi}
\frac{\partial\pi}{\partial g}
\frac{\partial g}{\partial\theta}.
\end{equation}
Thus, the reward itself is not differentiated through CityLearn. It determines
$\hat A_{it}$, which tells PPO whether to increase or decrease the probability of
the stored action; automatic differentiation then converts this probability
error into a gradient for the router.

For example, suppose two frozen experts assign probabilities 0.8 and 0.2 to the
executed action, and the router initially gives them weights 0.5 and 0.5. The
mixture assigns probability $0.5(0.8)+0.5(0.2)=0.5$ to that action. If its
advantage is positive, increasing the first router weight to 0.6 and decreasing
the second to 0.4 raises the mixture probability to
$0.6(0.8)+0.4(0.2)=0.56$, without changing either expert. This is the source of
the policy loss and router gradient during the frozen-expert stage.

The implemented actor objective also includes exploration and routing-balance
terms,
\begin{align}
\mathcal{L}_{\mathrm{actor}}
&=\mathcal{L}_{\mathrm{PPO}}-c_H\mathcal{H}(\pi)
+c_B\mathcal{L}_{\mathrm{balance}},\\
\mathcal{L}_{\mathrm{balance}}
&=\frac{1}{K}\sum_{k=1}^{K}(\bar g_k-\bar p_k)^2,
\end{align}
where $\bar g_k$ is mean expert usage in the batch and $\bar p_k$ is the original
group proportion. The balance term discourages routing collapse but is not an
additional CityLearn reward. The critic remains trainable and continues to
improve the value estimate even while the expert actors are frozen.

Several safeguards work together to reduce routing collapse. The router starts
with nearly equal expert probabilities, and training samples experts from these
probabilities instead of always taking the largest one. The entropy term
discourages the probabilities from becoming too concentrated, the original-group
prior and gradual increase of $\alpha$ preserve early diversity, the balance loss
discourages one expert from receiving nearly all buildings, and PPO clipping
limits abrupt probability changes. These measures reduce, but do not completely
eliminate, the possibility that an expert is rarely selected.

Freezing improves stability because it removes a moving-target problem. If the
router and every expert change rapidly at the same time, the router is trying to
learn which expert is suitable while the meaning and behaviour of every expert
are also changing. With pretrained experts fixed, their action distributions are
reliable reference points, and the smaller router only has to learn how to combine
them. Together with the safeguards above, this staged freezing makes learning
more stable. In the stable full-expert experiment, episodes 1--500 update the
router while freezing all experts;
episodes 501--1000 freeze the learned router and update only the small action
heads at a low learning rate. Keeping the encoders and TarMAC module fixed
preserves their pretrained latent spaces and communication behaviour. During
deterministic testing, each building directly uses the expert with the largest
routing probability.


\section{Router Training Sequence}

The soft-router experiments evolved through four designs:
\begin{enumerate}
  \item \textbf{End-to-end legacy training.} Router and actors were trained
  together from randomly initialized policies. Temperature, warm-up duration,
  and final prior weight were varied. These runs established that a router could
  assign different selection probabilities to the experts, but policy learning
  was unstable.
  \item \textbf{Pretrained two-stage training.} A fixed-group checkpoint using
  the strong 3fA feature combination identified earlier was loaded first.
  Experts were frozen for 200 episodes before
  limited joint fine-tuning, with temperature 0.5 or 0.7, final dynamic weight
  1.0 or 0.7, and with/without the seven capacity-aware router inputs.
  \item \textbf{Three-stage training.} Experiments separated static actor
  training, router-only learning, and dynamic actor adaptation. Both a shared-
  encoder version and a full-expert version were tested. The shared version
  attempted to put every head in one latent space; the full-expert version
  retained matched encoder--communication--head branches.
  \item \textbf{Stable two-stage enhancement.} The best fixed 3fA checkpoint
  was treated as an anchor. The most conservative run trained only the router
  for 500 episodes (\texttt{routeronly500}). The stronger stable-head family
  then allowed only action heads to adapt at $10^{-5}$ while encoders and TarMAC
  remained frozen and the router was frozen during the adaptation phase.
\end{enumerate}

This sequence tests the central question directly. If soft routing could replace
clustering, the end-to-end and shared three-stage runs should learn good experts
without a strong prior. If it is instead an enhancement, the pretrained and
stable variants should be better because they preserve a useful initial
partition.


\section{Completed Experiments}

At the current stage, the following experiments have been completed:
\begin{itemize}
  \item reward-weight comparisons using the baseline, stricter-comfort,
  controlled degree-hours, and load-tracking-focused parameter settings;
  \item comparison of the original TarMAC communication with TarMAC Hybrid, as
  well as the ReLU, Linear, and Gated fusion ablations;
  \item grouping-method experiments comparing K-means, GMM, agglomerative
  clustering, and balanced spectral clustering under the same 5F setting;
  \item grouping-feature experiments covering the full 9F control profile, the
  5F, 4F, and 3F feature-count ablations, and the nine representative 3F
  combinations A--I with multiple random seeds;
  \item multi-seed stability and cross-climate experiments for the selected 3fA
  feature roles in Vermont and Texas;
  \item Soft Router experiments that gradually improved the training process from
  an unstable initial design to the final stable version. The final version first
  keeps the pretrained actors unchanged while training the selector, and then
  keeps the learned selector fixed while making small adjustments to the actor
  output layers.
\end{itemize}


\section{Current Limitations}

The final stable Soft Router can be trained reliably and provides a useful
performance enhancement over fixed grouping. Its main limitation is that this
stability depends on strong pretrained actors produced by an earlier fixed
grouping; it still cannot learn equally strong actors and routing from random
initialization. Therefore, Soft Router currently improves an existing grouping
rather than replacing clustering completely. Validation on more climates and
portfolio sizes is still needed.
